\subsection{\texttt{TPARTMIN}}
\noindent\textbf{Bundle encoding.} UOP entry 38, opcode \texttt{0x1043}. Bundle: \texttt{B.OP + B.ARG + B.ATTR + B.IOTI}. \texttt{B.OP.Opcode}=\texttt{0x1043}. \texttt{B.IOTI}: selector=\texttt{111 last}, \texttt{SrcTile0}=\texttt{src0}, \texttt{SrcTile1}=\texttt{src1}, \texttt{DstTile}=\texttt{dst}, \texttt{SizeCode}=profile tile size. \texttt{B.ATTR} carries partial-valid-region policy.\par\smallskip
{\small
\paragraph{PTO ISA description.}
Partial elementwise min with implementation-defined handling of mismatched valid regions.
\paragraph{Example.}
Same pattern as \texttt{TPARTADD} $+$ but with min instead of add.
\par\smallskip
\paragraph{PTO ISA operation.}
For each element \texttt{(i, j)} in the destination valid region:

\[
\mathrm{dst}_{i,j} =
\begin{cases}
\min(\mathrm{src0}_{i,j}, \mathrm{src1}_{i,j}) & \text{if both inputs are defined at } (i,j) \\
\mathrm{src0}_{i,j} & \text{if only src0 is defined at } (i,j) \\
\mathrm{src1}_{i,j} & \text{if only src1 is defined at } (i,j)
\end{cases}
\]
}\par\smallskip
\paragraph{Notes.}
Same NaN handling as \texttt{TPARTMAX}.
\paragraph{Microcode site table.}
{\scriptsize
\begin{longtable}{@{}R{0.055\linewidth}L{0.23\linewidth}L{0.31\linewidth}L{0.22\linewidth}@{}}
\toprule
Seq & Micro instruction & WO[63:32] fields & WM[31:0] fields \\
\midrule
0 & \texttt{u\_vec.vbr} & \texttt{opc=0x0B, x=0, fl=alu, seq=0, kind=c} & \texttt{entry=38, cnt=0, res=vec, var=0} \\
1 & \texttt{u\_vec.vsts} & \texttt{opc=0x2A, x=0, fl=mem, seq=1, kind=b} & \texttt{entry=38, cnt=2, res=vec, var=0} \\
2 & \texttt{u\_vec.mem\_bar} & \texttt{opc=0x00, x=0, fl=setup, seq=2, kind=u} & \texttt{entry=38, cnt=1, res=vec, var=0} \\
3 & \texttt{u\_vec.vlds} & \texttt{opc=0x17, x=0, fl=mem, seq=3, kind=b} & \texttt{entry=38, cnt=2, res=vec, var=0} \\
4 & \texttt{u\_vec.vsts} & \texttt{opc=0x2A, x=0, fl=mem, seq=4, kind=b} & \texttt{entry=38, cnt=2, res=vec, var=0} \\
5 & \texttt{u\_vec.mem\_bar} & \texttt{opc=0x00, x=0, fl=setup, seq=5, kind=u} & \texttt{entry=38, cnt=1, res=vec, var=0} \\
6 & \texttt{u\_vec.vlds} & \texttt{opc=0x17, x=0, fl=mem, seq=6, kind=b} & \texttt{entry=38, cnt=2, res=vec, var=0} \\
7 & \texttt{u\_vec.vlds} & \texttt{opc=0x17, x=0, fl=mem, seq=7, kind=b} & \texttt{entry=38, cnt=2, res=vec, var=0} \\
8 & \texttt{u\_vec.vmin} & \texttt{opc=0x1C, x=0, fl=alu, seq=8, kind=b} & \texttt{entry=38, cnt=2, res=vec, var=0} \\
9 & \texttt{u\_vec.vsts} & \texttt{opc=0x2A, x=0, fl=mem, seq=9, kind=b} & \texttt{entry=38, cnt=2, res=vec, var=0} \\
\bottomrule
\end{longtable}
}
\paragraph{Microcode body DSL.}
\begin{lstlisting}[style=ptocode]
def uop_tpartmin(src0, src1, dst):
    dtype = dst.element_type
    valid_rows, valid_cols = dst.valid_shape
    src0_valid_rows, src0_valid_cols = src0.valid_shape
    src1_valid_rows, src1_valid_cols = src1.valid_shape
    lanes = pto.get_lanes(dtype)

    pad_scalar = pto.PadValue.MAX.eval(dtype)
    pad_vec = pto.vbr(pad_scalar)
    for row in range(0, valid_rows, 1):
        remained = valid_cols
        for col in range(0, valid_cols, lanes):
            mask, remained = pto.make_mask(dtype, remained)
            pto.vsts(pad_vec, dst[row, col:], mask)

    pto.mem_bar(pto.BarrierType.VST_VLD)

    for row in range(0, src0_valid_rows, 1):
        remained = src0_valid_cols
        for col in range(0, src0_valid_cols, lanes):
            mask, remained = pto.make_mask(dtype, remained)
            vec0 = pto.vlds(src0[row, col:])
            pto.vsts(vec0, dst[row, col:], mask)

    pto.mem_bar(pto.BarrierType.VST_VLD)

    for row in range(0, src1_valid_rows, 1):
        remained = src1_valid_cols
        for col in range(0, src1_valid_cols, lanes):
            mask, remained = pto.make_mask(dtype, remained)
            vec_dst = pto.vlds(dst[row, col:])
            vec1 = pto.vlds(src1[row, col:])
            result = pto.vmin(vec_dst, vec1, mask)
            pto.vsts(result, dst[row, col:], mask)

    return
\end{lstlisting}
\Needspace{0.34\textheight}
